Solve an exercise for determining the optimal contract with symmetric information

An entrepreneur has the opportunity to invest \( 100 \) monetary units (MU) in a technology generating the following
returns after one year: With a probability of \( p_1 = 0.95 \)  the technology generates a revenue of \( y_1 = 150 \text{MU} \), with
probability \( p_2 = 0.03 \) the revenue is \( y_2 = 100 \text{MU} \), and with probability \( p_3 = 0.02 \) the revenue is \( y_3 = 50 \text{MU} \). 
The entrepreneur does not have any equity. A bank without market power offers a loan to the entrepreneur. As an
alternative, the bank has the opportunity to invest its funds at the interbank market at a riskless interest rate of 10\%. 
Calculate the nominal repayment value that the bank will claim from the entrepreneur, if the bank’s utility function is given by \( \text{U}_L(x) = 0.75x \) and the entrepreneur’s utility is \( \text{U}_B(x) = 0.5x \).
How can the bank’s and entrepreneur’s risk preferences be described?

The optimization problem is
\( \max_R \mathbb{E}[U_B(y)] = \mathbb{E}[\frac{1}{2}\sqrt{y - R(y)}] \) s.t. \( \mathbb{E}[U_L(R(y))] > U^0_L \).
It holds \( U^0_L = 0.75\sqrt{100\cdot 1.1} \). Denote \( R = R(y) \), then the constraint reads 
\( \mathbb{E}[0.75\sqrt{R}] \geq 0.75\sqrt{100\cdot 1.1} \)
Solving this for \( R \) with given probabilities yields \( R \geq 111.8391 \), therefore the borrowers utility
is
\( 0.95 \cdot 0.5 \cdot \sqrt{150 - 111.84} + 0.03\cdot 0 + 0.02 \cdot 0 = 2.934 \).
Note the constraint \( y \geq R \) for the borrower (as otherways they wouldn't undertake the project as the cost is higher than the return).
To analyze the their risk preferences, one does the usual first and second derivative of the \( U \)'s and checks convexity/concavity.


Solve an exercise for determining the optimal contract with asymmetric information

A borrower wants to finance her investment project exclusively via credit in the amount of 100 MU. If the project
is successful, the profit will be 20, in case it is unsuccessful, there will be a loss of 20. The risk neutral lender
knows that there are two types of borrowers - “safe” and “risky” ones. The safe borrowers’ probability of success
is \( p_S = 1 \); however, the risky borrowers’ probability of success is only \( p_R = 0.9 \). Furthermore, the share of safe borrowers
within the population is known \( \pi = 0.6 \) of all borrowers are safe and \( 1 - \pi \) are risky.
Calculate the interest rate the lender will claim if he requires an expected interest of \( r_f = 0.04 \).

Lenders constraint
\( \pi(p_S\cdot R + (1-p_S)\cdot 80) + (1 - \pi)(p_r \cdot R + (1 - p_R)\cdot 80) = 100\cdot 1.04 \)
Solving this for \( R \) leads to 
\( R = 105 \) so \( r = \frac{R}{100} - 1 = 0.05 \).
The required interest is higher than the alternative investment because it contains a risk premium that accounts for default risk.



What is the solution for the option price in the binomial tree model? 
What is the general approach? How to hedge a claim?

Let the initial value of the asset be \( S_0 \), the risk-free interest rate \( r \) constant and for each \( t \in \{0, \dots, T\} \)
\( S_t = S_0u^{i}d^{t-i} \) for some \( i \leq t \). A general claim \( C_t = f(S_t, t) \) can then be evaluated recursively
by starting from \( S_T \) via \( C_t = e^{-rT}(\frac{e^{rT} - d}{u - d}C_{t+1,u} + \frac{u - e^{rT}}{u - d}C_{t+1,d} \)
where \( C_{t+1,u} \) is the value of the claim at \( t+1 \) for an up-move of \( S_t \).

The risk-neutral probability is \( q = \frac{e^{rT} - d}{u - d} \), thus an arbitrage free market satisfies \( d < e^{rT} < u \).

The hedging portfolio with positions \( (\delta_t,B_t) \) in the asset and some savings account with rate \( r \)
is \( C_t = S_t\delta_t + B_te^{rT}\). The solution is then
\( \delta_t = \frac{C_{t+1,u} - C_{t+1,d}}{S_tu - S_td} \)
\( B_t = \frac{C_{t+1,u} - \delta_t S_t u}{e^{rT}} \)

under the assumptions of the Cox-Ross-Rubinstein model \( CRR \), that is \( u = e^{\sigma\sqrt{\frac{T}{N}}},\ d = \frac{1}{u} \),
where \( T \) is the overall time and \( N \) the number of steps, the prices are recombinant, meaning \( S_t ud = S_t \) and
the price for a call simplifies to
\( C_t e^{-rT} \sum_{i=0}^N \binom{N}{i}q^i(1 - q)^{N-i}\max\{Su^id^{N-i} - K, 0\}\).



What should you know in regards to the Black-Scholes-Merton model?

Mostly the same as for advanced derivatives, therefore no extra anki cards are given.




What is a coherent risk measure?

Monetary Risk Measure:
A mapping \(\rho : \mathcal{X} \to \mathbb{R}\) is called a monetary risk measure if it satisfies the following properties:  
1. Inverse monotonicity: If \(\mathcal{X}(\omega) \leq \mathcal{Y}(\omega)\) for all \(\omega \in \Omega\), then 
   \(\rho(\mathcal{X}) \geq \rho(\mathcal{Y})\).  
2. Cash invariance: For any \(m \in \mathbb{R}\), \(\rho(\mathcal{X} + m) = \rho(\mathcal{X}) - m\).

Interpretation:  
1. Property 1 states that the risk of a position \(\mathcal{Y}\) is smaller than the risk of a position \(\mathcal{X}\), 
   if the future value of \(\mathcal{Y}\) is at least as large as \(\mathcal{X}\) for any scenario.  
2. Property 2 states that risk is measured on a monetary scale: if \(m\) are added to \(\mathcal{X}\), 
   then the risk of \(\mathcal{X}\) is exactly reduced by this amount.

A monetary risk measure \(\rho\) on \(\mathcal{X}\) is called a convex risk measure 
if it is a convex functional and thus does not penalize diversification.
A convex risk measure is called coherent if it satisfies in addition positive homogeneity, i.e.,
\(
\rho(\lambda X) = \lambda \rho(X) \quad \text{for all } X \in \mathcal{X} \text{ and } \lambda \geq 0.
\)

Remarks:
Each coherent risk measure \(\rho\) is sub-additive, i.e.,
\(
\rho(X + Y) \leq \rho(X) + \rho(Y) \quad \text{for all } X, Y \in \mathcal{X}.
\)
Positive Homogeneity is a scaling property.
- Note that this property might be inappropriate for risk measurement, due to concentration and liquidity risk.
- This criticism motivated the notion of convex risk measures, introduced by Föllmer and Schied (2002).


State some of the relevant risk measures in financial intermediation

Value at Risk (not coherent, only monetary)
\( \text{VaR}_{\alpha}(Y) = \inf\{y|F_Y(y) \geq \alpha\} = F^{-1}(\alpha) \)
if the inverse exists.

Economic capital (EC)
is defined as the difference between the value-at-risk and the expected loss.
\( \text{EC}_\alpha(L) = \text{VaR}_{\alpha}(L) - \mathbb{E}[L] \)

The Conditional Value-at-Risk \( \text{CVaR}_{\alpha}(Y) \) of confidence level
\( \alpha \in (0,1) \) quantifies the expected loss incurred in the \( (1 - \alpha)100 \)% worst cases.
\( \text{CVaR}_{\alpha}(Y) = \int_{-\infty}^\infty z dG^\alpha_Y(z) \)
where 
\[ G^\alpha_Y = \begin{cases}
    0, & z < \text{VaR}_\alpha(Y), \\
    \frac{F_Y(z) - a}{1 - a},& \text{ else}
\end{cases}
\]

This usually simplifies to
\( \text{CVaR}_{\alpha}(Y) = \frac{1}{\alpha}\int_0^\alpha \text{VaR}_{\beta}(Y) d\beta \)



Give an example how to determine the \( \text{CVaR}_{\alpha}^\pm \)

Upper and Lower \( \text{CVaR}_{\alpha}^\pm \)
\( \text{CVaR}_{\alpha}^{+} = \mathbb{E}[Y|Y > \text{VaR}_{\alpha}(Y) \)
\( \text{CVaR}_{\alpha}^{-} = \mathbb{E}[Y|Y \geq \text{VaR}_{\alpha}(Y) \)

If \( F_Y(\text{VaR}_\alpha(Y)) < 1 \), it holds 
\( \text{CVaR}_{\alpha}(Y) = \lambda_\alpha(Y)\text{VaR}_\alpha(Y) + (1 - \lambda_\alpha)\text{CVaR}_{\alpha}^{+}\)
where \( \lambda_\alpha = \frac{F_Y(\text{VaR}_\alpha(Y)) - \alpha}{1 - \alpha} \),
else \( \text{VaR}_\alpha(Y) = \text{CVaR}_{\alpha}(Y) \).

INSERT IMAGE 5e, 22 
INSERT IMAGE 5e, 24

Compare different risk measures qualitatively

Expected Loss:
(+) Information on average portfolio loss; direct connection to provisions; Coherent
(-) No information on the shape of the loss distribution.
Loss standard deviation:
(+) Information on loss uncertainty / risk and scale of the loss distribution.
(-) Less information for asymmetric distribution; not coherent.
Value-at-Risk:
(+) Intuitive and commonly used; confidence level interpretation; actively used in
banks; capital calculations and risk-adjusted performance measures
(-) No information on shape, only info on one percentile; difficult to compute and
interpret at very high percentiles; not coherent (lacks subadditivity).
Economic capital:
(+) Intuitive and commonly used; confidence level interpretation; actively
used in banks; capital calculations and risk-adjusted performance measures.
(-) No information on shape, only info on one percentile; difficult to compute
and interpret at very high percentiles; not coherent
Conditional Value-at-Risk:
(+) Coherent measure of risk at a given confidence level, increasingly
popular in banks.
(-) Less intuitive than VaR; only tail and distribution information for the
given percentile; computational issues.
